\documentclass[multi=tikzpicture,border=9pt]{standalone}
\usepackage[UTF8]{ctex}
\usepackage{amsmath,amssymb}
\usepackage{tikz}
\usetikzlibrary{arrows.meta,calc,patterns,decorations.pathmorphing}
\definecolor{main}{HTML}{17365D}
\definecolor{accent}{HTML}{AC3931}
\tikzset{every picture/.style={font=\small,line width=.75pt},
>={Stealth[length=2mm]},arr/.style={->},
box/.style={draw=main,rounded corners=2pt,fill=main!4,text width=3.6cm,minimum height=1cm,align=center},
note/.style={align=center,text width=9cm}}
\begin{document}

\begin{tikzpicture}
\draw[main,thick] (0,2) circle (2);\fill (0,2) circle (2pt);\node[above right] at (0,2){$O$};\draw[accent,very thick] (0,2)--(-2,2);\draw[main,thick] (-3,-2)--(3,-2);\draw[arr,accent] (2,-2)--(.5,-2) node[midway,below]{$I$};\draw[<->] (2.5,-2)--(2.5,0) node[midway,right]{$R$};\draw[<->] (-2.7,-2)--(-2.7,2) node[midway,left]{$2R$};\draw[arr] (-.6,4.5) arc (110:70:2);\node[above] at (0,4.5){顺时针 $\omega$};\node[note] at (0,-3.4){导线到轮缘最低点间隙为 $R$，到轮心为 $2R$};
\end{tikzpicture}

\begin{tikzpicture}
\draw[main,thick] (-3,0)--(3,0);\draw[arr,accent] (1,0)--(-1,0) node[midway,below]{$I$};\foreach \x in {-2,0,2}{\foreach \y in {1,2,3}{\node at (\x,\y){$\otimes$};}}\node[note] at (0,-1.5){导线上方磁场向纸内；$B=\mu_0I/(2\pi d)$};
\end{tikzpicture}

\begin{tikzpicture}
\draw[main,very thick] (0,2)--(4,2);\draw[main] (-1,0)--(5,0);\node[below] at (0,2){$O$};\node[below] at (4,2){$P$};\draw[<->] (-.7,0)--(-.7,2) node[midway,left]{$2R$};\draw[arr,accent] (2,2)--(2,1) node[right]{$\mathbf v$};\node at (3,1){$\otimes\ \mathbf B$};\draw[arr] (2,2.3)--(3.5,2.3) node[above]{$\mathbf v\times\mathbf B$};\node[note] at (2,-1.5){磁场沿水平辐条恒定，速度随半径线性增加};
\end{tikzpicture}

\begin{tikzpicture}
\draw[main,very thick] (0,2)--(0,4);\draw[main] (-3,0)--(3,0);\node[right] at (0,2){$O$};\node[right] at (0,4){$P$};\draw[<->] (-.8,0)--(-.8,2) node[midway,left]{$2R$};\draw[<->] (-.8,2)--(-.8,4) node[midway,left]{$R$};\draw[arr,accent] (0,3)--(1,3) node[above]{$\mathbf v$};\draw[arr] (.3,2.8)--(.3,3.8);\node[note] at (0,-1.4){向上：导线距离为 $2R+r$，动生积分含自然对数};
\end{tikzpicture}

\begin{tikzpicture}
\draw[arr] (0,0)--(6,0);\foreach \x/\h/\lab in {1/1.89/向上,3/2.5/水平,5/3.86/向下}{\fill[main!20] (\x-.4,0) rectangle (\x+.4,\h);\draw[main] (\x-.4,0) rectangle (\x+.4,\h);\node[below] at (\x,0){\lab};}\node[note] at (3,-1.5){共同系数 $\mu_0I\omega R/(2\pi)$\\系数依次为 $0.189070$、$0.25$、$0.386294$};
\end{tikzpicture}
\end{document}
